\section*{Source: Demand Letter.docx}
\textbf{Category}: source-8
\textbf{Page}: 1

\subsection*{Demand for Adjudication Under the Mandamus Act and Administrative Procedure Act – Jessica Marie Thompson}

[Attorney’s Name]
[Law Firm Name]
[Street Address]
[City, State, ZIP Code]
[Phone Number]
[Email Address]

Date: 2026-06-12

Employer Contact:
- Name: [Employer’s Name]
- Title: [Title]
- Company: [Company Name]
- Address: [Street Address], City, State, ZIP Code

RE: Demand for Adjudication under the Mandamus Act – Jessica Marie Thompson

\subsection*{Dear [Employer’s Name],}

Introduction \& Jurisdiction
This letter is submitted by [Law Firm Name] on behalf of [Employer Name] (the “Petitioner”) in support of its EB-1A I-140 and I-485 petition for Jessica Marie Thompson (the “Beneficiary”). In this submission, we outline the applicability of mandamus claims based upon the established procedural framework. The petition history currently notes an I-140 filing on [Date] with Receipt No.: [Number] and an I-485 filing on [Date] with Receipt No.: [Number]. In addition, an RFE was issued on [Date] with a response filed on [Date]. This demand is made under 28 U.S.C. § 1361 and 5 U.S.C. § 555(b) due to the significant delay beyond the customary 60-day adjudication period.

Factual Background
The factual background presented details the petitioner’s status as well as the Beneficiary’s extensive qualifications as documented in the accompanying degree certificates, passport and I-94, and visa pages. Jessica Marie Thompson, born on 08/14/1993, holds a Bachelor of Technology in Computer Science and Engineering from Veermata Jijabai Technological Institute with First Class with Distinction, establishing her unique qualifications and specialized expertise in her field. Although detailed job descriptions and employer specifics are not provided, the evidence clearly demonstrates the critical need for the Beneficiary’s advanced skills. Thorough documentation supports the timeline and qualifications relevant to the EB-1A petition, reinforcing the indisputable right of the petition.

Legal Standard for Mandamus
The legal standard for mandamus requires proving a clear right to a final decision on the pending petitions, a non-discretionary duty by USCIS to act within a reasonable time frame, and the absence of any adequate alternative means of relief. The petitioner’s entitlement is clearly established by the submitted filings and extensive supporting documents. USCIS’s failure to adjudicate within their established guidelines creates an environment of undue delay, thereby satisfying the requirements under 28 U.S.C. § 1361 and 5 U.S.C. § 555(b). This section elaborates on the direct application of these legal standards as they relate to the current case, warranting prompt judicial intervention.

Demand for Relief
In light of the extended delay and the clear statutory requirements, we demand the immediate adjudication of both the EB-1A I-140 and I-485 petitions within 14 days from receipt of this letter. This demand is based on the firm statutory authority provided by 28 U.S.C. § 1361 and 5 U.S.C. § 555(b). Failure to resolve this matter promptly will result in significant consequences for both the petitioner and the Beneficiary, including operational delays and financial losses. It is imperative that USCIS provide a final decision without further delay to prevent additional prejudice to both employer operations and the legal status of the Beneficiary.

Prejudice \& Hardship
The economic and operational impact on the petitioner, alongside severe personal and professional detriment to the Beneficiary, demands immediate resolution. Prolonged delays risk project interruptions, financial setbacks estimated at approximately \$[Amount] per week, and potential contractual breaches. Moreover, any delay in adjudication jeopardizes the Beneficiary’s work authorization, thereby risking her lawful stay and causing familial and dependents’ instability. The significant hardships detailed herein underscore the inadequacy of monetary relief alone, demonstrating that only prompt judicial intervention via mandamus can fully remediate this circumstance and forestall irreparable injury.

Conclusion \& Next Steps
Based on the foregoing, we request that USCIS issue a final decision on the pending petitions no later than 14 days from the date of service of this letter. Please confirm receipt by emailing [Attorney’s Email] or faxing [Fax Number]. Absent timely action, we are prepared to file a Writ of Mandamus in the U.S. District Court for the District of [District] and seek reimbursement for EAJA fees and associated costs. [Attorney’s Name] remains available to provide further details or documentation as needed to facilitate an expedited and definitive resolution to this matter.

very truly yours,
\_\_\_\_\_\_\_\_\_\_\_,
[Attorney’s Full Name], Esq.
[Law Firm Name]